% !TEX program = xelatex
% =====================================================================
%  PhotoLib 部长使用指南
%  编译：xelatex 部长使用指南.tex（连续两次以生成目录）
% =====================================================================
\documentclass[cn,11pt,green]{elegantbook}

% =====================================================================
%  elegant-style.tex —— PhotoLib 使用指南三份文档共用的样式与宏定义
%  由 管理员/部长/校区负责人 三份指南通过 % =====================================================================
%  elegant-style.tex —— PhotoLib 使用指南三份文档共用的样式与宏定义
%  由 管理员/部长/校区负责人 三份指南通过 % =====================================================================
%  elegant-style.tex —— PhotoLib 使用指南三份文档共用的样式与宏定义
%  由 管理员/部长/校区负责人 三份指南通过 \input{elegant-style} 引入。
%  依赖：ElegantBook 文档类（须以 XeLaTeX 编译）。
% =====================================================================

% ---- 颜色（与 ElegantBook green 主题协调）----
\definecolor{plMain}{RGB}{0,120,80}      % 主色：墨绿
\definecolor{plInk}{RGB}{38,50,56}       % 正文强调
\definecolor{admTipBg}{RGB}{236,247,240}
\definecolor{admTipFr}{RGB}{34,139,88}
\definecolor{admNoteBg}{RGB}{232,240,254}
\definecolor{admNoteFr}{RGB}{25,103,210}
\definecolor{admWarnBg}{RGB}{253,237,237}
\definecolor{admWarnFr}{RGB}{198,40,40}
\definecolor{admStepBg}{RGB}{245,243,255}
\definecolor{admStepFr}{RGB}{103,58,183}

% ---- tcolorbox（ElegantBook 已载入，这里补声明所需库）----
\usepackage{tcolorbox}
\tcbuselibrary{skins,breakable}

% 统一的四种“提示框”环境，均支持可选自定义标题：
%   \begin{tipbox}[自定义标题] ... \end{tipbox}
\newtcolorbox{tipbox}[1][小贴士]{%
  enhanced, breakable,
  colback=admTipBg, colframe=admTipFr, coltitle=white,
  fonttitle=\bfseries\sffamily, title={#1},
  boxrule=0.6pt, arc=3pt, left=8pt, right=8pt, top=6pt, bottom=6pt}

\newtcolorbox{notebox}[1][说明]{%
  enhanced, breakable,
  colback=admNoteBg, colframe=admNoteFr, coltitle=white,
  fonttitle=\bfseries\sffamily, title={#1},
  boxrule=0.6pt, arc=3pt, left=8pt, right=8pt, top=6pt, bottom=6pt}

\newtcolorbox{warnbox}[1][注意]{%
  enhanced, breakable,
  colback=admWarnBg, colframe=admWarnFr, coltitle=white,
  fonttitle=\bfseries\sffamily, title={#1},
  boxrule=0.6pt, arc=3pt, left=8pt, right=8pt, top=6pt, bottom=6pt}

\newtcolorbox{stepbox}[1][操作步骤]{%
  enhanced, breakable,
  colback=admStepBg, colframe=admStepFr, coltitle=white,
  fonttitle=\bfseries\sffamily, title={#1},
  boxrule=0.6pt, arc=3pt, left=8pt, right=8pt, top=6pt, bottom=6pt}

% ---- 表格 ----
\usepackage{booktabs}
\usepackage{array}
\usepackage{longtable}
\usepackage{makecell}
\renewcommand{\arraystretch}{1.28}

% 左对齐、可换行的定宽列
\newcolumntype{L}[1]{>{\raggedright\arraybackslash}p{#1}}

% ---- 行内“键位/菜单/字段”标记 ----
\usepackage{relsize}
% 菜单路径：系统管理 \menu{数据备份}
\newcommand{\menu}[1]{{\sffamily\bfseries\color{plMain}#1}}
% 界面按钮/字段名
\newcommand{\ui}[1]{\textsf{「#1」}}
% 权限码等等宽标识；在每个下划线后允许换行，避免长标识顶出表格边界
\newcommand{\code}[1]{{\ttfamily\smaller[0.5]%
  \renewcommand{\_}{\textunderscore\penalty0\hspace{0pt}}#1}}
% 角色徽标
\newcommand{\role}[1]{{\sffamily\bfseries\color{admNoteFr}#1}}

% ---- 超链接 ----
\hypersetup{
  colorlinks=true,
  linkcolor=plMain,
  urlcolor=admNoteFr,
  citecolor=plMain}
 引入。
%  依赖：ElegantBook 文档类（须以 XeLaTeX 编译）。
% =====================================================================

% ---- 颜色（与 ElegantBook green 主题协调）----
\definecolor{plMain}{RGB}{0,120,80}      % 主色：墨绿
\definecolor{plInk}{RGB}{38,50,56}       % 正文强调
\definecolor{admTipBg}{RGB}{236,247,240}
\definecolor{admTipFr}{RGB}{34,139,88}
\definecolor{admNoteBg}{RGB}{232,240,254}
\definecolor{admNoteFr}{RGB}{25,103,210}
\definecolor{admWarnBg}{RGB}{253,237,237}
\definecolor{admWarnFr}{RGB}{198,40,40}
\definecolor{admStepBg}{RGB}{245,243,255}
\definecolor{admStepFr}{RGB}{103,58,183}

% ---- tcolorbox（ElegantBook 已载入，这里补声明所需库）----
\usepackage{tcolorbox}
\tcbuselibrary{skins,breakable}

% 统一的四种“提示框”环境，均支持可选自定义标题：
%   \begin{tipbox}[自定义标题] ... \end{tipbox}
\newtcolorbox{tipbox}[1][小贴士]{%
  enhanced, breakable,
  colback=admTipBg, colframe=admTipFr, coltitle=white,
  fonttitle=\bfseries\sffamily, title={#1},
  boxrule=0.6pt, arc=3pt, left=8pt, right=8pt, top=6pt, bottom=6pt}

\newtcolorbox{notebox}[1][说明]{%
  enhanced, breakable,
  colback=admNoteBg, colframe=admNoteFr, coltitle=white,
  fonttitle=\bfseries\sffamily, title={#1},
  boxrule=0.6pt, arc=3pt, left=8pt, right=8pt, top=6pt, bottom=6pt}

\newtcolorbox{warnbox}[1][注意]{%
  enhanced, breakable,
  colback=admWarnBg, colframe=admWarnFr, coltitle=white,
  fonttitle=\bfseries\sffamily, title={#1},
  boxrule=0.6pt, arc=3pt, left=8pt, right=8pt, top=6pt, bottom=6pt}

\newtcolorbox{stepbox}[1][操作步骤]{%
  enhanced, breakable,
  colback=admStepBg, colframe=admStepFr, coltitle=white,
  fonttitle=\bfseries\sffamily, title={#1},
  boxrule=0.6pt, arc=3pt, left=8pt, right=8pt, top=6pt, bottom=6pt}

% ---- 表格 ----
\usepackage{booktabs}
\usepackage{array}
\usepackage{longtable}
\usepackage{makecell}
\renewcommand{\arraystretch}{1.28}

% 左对齐、可换行的定宽列
\newcolumntype{L}[1]{>{\raggedright\arraybackslash}p{#1}}

% ---- 行内“键位/菜单/字段”标记 ----
\usepackage{relsize}
% 菜单路径：系统管理 \menu{数据备份}
\newcommand{\menu}[1]{{\sffamily\bfseries\color{plMain}#1}}
% 界面按钮/字段名
\newcommand{\ui}[1]{\textsf{「#1」}}
% 权限码等等宽标识；在每个下划线后允许换行，避免长标识顶出表格边界
\newcommand{\code}[1]{{\ttfamily\smaller[0.5]%
  \renewcommand{\_}{\textunderscore\penalty0\hspace{0pt}}#1}}
% 角色徽标
\newcommand{\role}[1]{{\sffamily\bfseries\color{admNoteFr}#1}}

% ---- 超链接 ----
\hypersetup{
  colorlinks=true,
  linkcolor=plMain,
  urlcolor=admNoteFr,
  citecolor=plMain}
 引入。
%  依赖：ElegantBook 文档类（须以 XeLaTeX 编译）。
% =====================================================================

% ---- 颜色（与 ElegantBook green 主题协调）----
\definecolor{plMain}{RGB}{0,120,80}      % 主色：墨绿
\definecolor{plInk}{RGB}{38,50,56}       % 正文强调
\definecolor{admTipBg}{RGB}{236,247,240}
\definecolor{admTipFr}{RGB}{34,139,88}
\definecolor{admNoteBg}{RGB}{232,240,254}
\definecolor{admNoteFr}{RGB}{25,103,210}
\definecolor{admWarnBg}{RGB}{253,237,237}
\definecolor{admWarnFr}{RGB}{198,40,40}
\definecolor{admStepBg}{RGB}{245,243,255}
\definecolor{admStepFr}{RGB}{103,58,183}

% ---- tcolorbox（ElegantBook 已载入，这里补声明所需库）----
\usepackage{tcolorbox}
\tcbuselibrary{skins,breakable}

% 统一的四种“提示框”环境，均支持可选自定义标题：
%   \begin{tipbox}[自定义标题] ... \end{tipbox}
\newtcolorbox{tipbox}[1][小贴士]{%
  enhanced, breakable,
  colback=admTipBg, colframe=admTipFr, coltitle=white,
  fonttitle=\bfseries\sffamily, title={#1},
  boxrule=0.6pt, arc=3pt, left=8pt, right=8pt, top=6pt, bottom=6pt}

\newtcolorbox{notebox}[1][说明]{%
  enhanced, breakable,
  colback=admNoteBg, colframe=admNoteFr, coltitle=white,
  fonttitle=\bfseries\sffamily, title={#1},
  boxrule=0.6pt, arc=3pt, left=8pt, right=8pt, top=6pt, bottom=6pt}

\newtcolorbox{warnbox}[1][注意]{%
  enhanced, breakable,
  colback=admWarnBg, colframe=admWarnFr, coltitle=white,
  fonttitle=\bfseries\sffamily, title={#1},
  boxrule=0.6pt, arc=3pt, left=8pt, right=8pt, top=6pt, bottom=6pt}

\newtcolorbox{stepbox}[1][操作步骤]{%
  enhanced, breakable,
  colback=admStepBg, colframe=admStepFr, coltitle=white,
  fonttitle=\bfseries\sffamily, title={#1},
  boxrule=0.6pt, arc=3pt, left=8pt, right=8pt, top=6pt, bottom=6pt}

% ---- 表格 ----
\usepackage{booktabs}
\usepackage{array}
\usepackage{longtable}
\usepackage{makecell}
\renewcommand{\arraystretch}{1.28}

% 左对齐、可换行的定宽列
\newcolumntype{L}[1]{>{\raggedright\arraybackslash}p{#1}}

% ---- 行内“键位/菜单/字段”标记 ----
\usepackage{relsize}
% 菜单路径：系统管理 \menu{数据备份}
\newcommand{\menu}[1]{{\sffamily\bfseries\color{plMain}#1}}
% 界面按钮/字段名
\newcommand{\ui}[1]{\textsf{「#1」}}
% 权限码等等宽标识；在每个下划线后允许换行，避免长标识顶出表格边界
\newcommand{\code}[1]{{\ttfamily\smaller[0.5]%
  \renewcommand{\_}{\textunderscore\penalty0\hspace{0pt}}#1}}
% 角色徽标
\newcommand{\role}[1]{{\sffamily\bfseries\color{admNoteFr}#1}}

% ---- 超链接 ----
\hypersetup{
  colorlinks=true,
  linkcolor=plMain,
  urlcolor=admNoteFr,
  citecolor=plMain}


\title{PhotoLib 部长使用指南}
\subtitle{立项 · 发需求 · 定采用 · 看数据}
\author{\href{mailto:i@liujr.uk}{PhotoLib项目组}}
\institute{校园摄影部图片工作站}
\version{1.0}
\date{2026 年 9 月}
\extrainfo{本指南面向摄影部正副部长（MINISTER），用最直白的话讲清楚\\你每天要用到的每一个功能，跟着做就能上手。}
\cover{cover.png}

\begin{document}

\maketitle
\frontmatter
\tableofcontents
\mainmatter

% =====================================================================
\chapter{开始之前}
% =====================================================================

\begin{introduction}
  \item 你在系统里能做什么
  \item 怎么登录、第一次要注意什么
  \item 界面大概长什么样
\end{introduction}

\section{你的角色}

你是摄影部的\textbf{部长}（系统里的角色叫 \role{MINISTER}）。简单说，你负责\textbf{「组织生产」}这条线：

\begin{itemize}
  \item 立选题、发图片需求，把活儿派给各校区。
  \item 从大家交上来的图片里\textbf{挑出被采用的图}（标记「被引」）。
  \item \textbf{确认}校区负责人填报的工时。
  \item \textbf{导出}每位成员的被引张数和工时数据。
  \item 还能发布成员招募、维护文档中心、发站内消息、评选好图精选。
\end{itemize}

\begin{notebox}[你不需要操心的事]
  账号创建、权限分配、系统配置、数据备份这些是\textbf{管理员}的事，你在菜单里也看不到「系统管理」。
  你只管用好自己这条业务线就行。
\end{notebox}

\section{登录与第一次改密}

\begin{stepbox}[第一次登录]
  \begin{enumerate}
    \item 打开管理员给你的网址，输入\textbf{用户名（或邮箱）+ 密码}。
    \item 第一次登录，系统会\textbf{强制你改密码}。设一个只有你知道的强密码。
    \item 改完自动进入\ui{工作台}首页。
  \end{enumerate}
\end{stepbox}

\begin{tipbox}[忘了密码怎么办]
  系统\textbf{没有自助找回}。忘记密码时，找管理员帮你重置即可。重置后你下次登录会再改一次密码。
\end{tipbox}

\section{界面导览}

登录后，左边是\textbf{菜单栏}，右上角是\textbf{消息铃铛}和\textbf{头像}。你常用的菜单大致是：

\begin{table}[htbp]
  \centering
  \caption{部长常用菜单一览}
  \begin{tabular}{L{2.8cm} L{9.4cm}}
    \toprule
    \textbf{菜单} & \textbf{用来做什么} \\
    \midrule
    工作台 & 首页概览。 \\
    选题项目 & 立项、编辑、发布、标记完成。 \\
    图片需求 & 给校区派活儿；确认或退回交付。 \\
    图片库 & 浏览、下载、收藏图片，标记被引。 \\
    好图精选 & 评选和展示优秀作品。 \\
    工时记录 & 确认、退回、导出工时。 \\
    数据统计 & 按成员导出被引张数和工时。 \\
    通讯录 & 查看各校区成员名单（只读）。 \\
    成员招募 & 设计报名表、发布招募、查看报名。 \\
    文档中心 & 编写和发布部门文档。 \\
    消息中心 & 查看站内通知。 \\
    \bottomrule
  \end{tabular}
\end{table}

% =====================================================================
\chapter{选题项目}
% =====================================================================

\begin{introduction}
  \item 怎么新建一个选题
  \item 说明里怎么配图、排版
  \item 项目做完了怎么收尾
\end{introduction}

\section{新建选题}

进入\menu{选题项目}，点\ui{新建}。填写项目标题、说明等信息后保存。说明支持 \textbf{Markdown 排版}
（可以加标题、列表、加粗），你也可以在说明里\textbf{插入图片}。

\begin{tipbox}[给说明配图]
  在说明编辑器里上传图片即可，系统只接受 JPEG / PNG / WebP，单张不超过 5 MiB。
  这些图片通过登录后才能访问的安全地址读取，不会外泄。
\end{tipbox}

\section{编辑与发布}

新建后的项目可以继续编辑。项目要处于\textbf{进行中（ACTIVE）}状态，才能往里发布图片需求。

\section{标记完成}

项目做完后，把它\textbf{标记为完成}。要注意：

\begin{warnbox}[完成会「锁定」项目]
  项目一旦完成，就\textbf{不能再发布新需求、也不能再新增采用记录}了。所以请确认所有图片都收齐、
  该采用的都采用了，再点完成。
\end{warnbox}

% =====================================================================
\chapter{图片需求}
% =====================================================================

\begin{introduction}
  \item 怎么把一个拍摄任务派给校区
  \item 怎么一次派给多个校区
  \item 交付上来后怎么确认或退回
\end{introduction}

\section{需求是什么}

「需求」就是你在一个项目下发布的\textbf{一条拍摄任务}。校区负责人看到后可以\textbf{接单}，
拍完把图\textbf{交付}上来。一条需求可以\textbf{多个人一起接、一起交}。

需求会经历这样几个状态：草稿 → 已发布 → 有人接受 → 已提交 → 已完成（或中途取消）。

\section{新建并发布一条需求}

\begin{stepbox}[发布单条需求]
  \begin{enumerate}
    \item 进入某个\textbf{进行中}的项目，或从\menu{图片需求}进入。
    \item 新建需求，填写说明、\textbf{截止时间}（不能早于现在）。
    \item 需求刚建好是\textbf{草稿}，只有你（或管理员）能编辑、发布、取消。
    \item 点\ui{发布}。发布后，\textbf{同校区已启用的校区负责人}会收到通知。
  \end{enumerate}
\end{stepbox}

\begin{notebox}[数量可以不填]
  现在新建需求\textbf{不强制填「需要几张」}。如果你填了，至少要填 1。
\end{notebox}

\section{一次派给多个校区}

如果同一个任务要好几个校区一起拍，用\textbf{批量发布}：一次最多选 \textbf{50 个校区}，系统会为\textbf{每个校区}
各建一条独立的、直接进入「已发布」的需求。

\begin{tipbox}[部分成功也没关系]
  批量发布时，即使某个校区因为故障没建成功，\textbf{其他校区照样成功}。界面会把失败的校区留出来，
  你可以单独重试，不用整批重来。
\end{tipbox}

\section{确认或退回交付}

校区把图交上来（需求变成「已提交」）后，你来\textbf{确认}或\textbf{退回}：

\begin{itemize}
  \item 觉得没问题，就\textbf{确认}。
  \item 觉得要重拍或补拍，就\textbf{退回}，让对方继续。
\end{itemize}

\begin{notebox}[允许零图片交付]
  系统现在允许一条需求「没有图片也能提交」。所以确认前请\textbf{自己核对交付内容}，别只看状态。
\end{notebox}

% =====================================================================
\chapter{图片库与采用}
% =====================================================================

\begin{introduction}
  \item 怎么找图、下载图、收藏图
  \item 「标记被引」是怎么回事
\end{introduction}

\section{浏览与下载}

进入\menu{图片库}，可以按关键词、状态等条件检索图片。你能看到全站的图片。

\begin{itemize}
  \item \textbf{下载}：单张下载，或勾选多张\textbf{批量打包下载}（一次最多 200 张，打成 ZIP）。
  \item \textbf{收藏}：点图片上的星标即可收藏，之后在\menu{收藏图片}里集中查看。
\end{itemize}

\begin{tipbox}[下载链接是临时的]
  图片下载用的是\textbf{短时有效的安全链接}（默认 15 分钟）。链接过期只要重新点一下即可，属于正常现象。
\end{tipbox}

\section{标记被引（采用）}

「被引」就是\textbf{这张图被正式采用了}。这是统计成员贡献的核心动作。

\begin{stepbox}[采用一张图]
  \begin{enumerate}
    \item 图片得\textbf{先在目标项目的相册里}，你才能把它标记为该项目的采用图。
    \item 在图片详情或图库里，把它\textbf{标记被引}到某个项目。
    \item 从图库批量采用，一次最多 200 张，且目标项目必须是\textbf{进行中}。
  \end{enumerate}
\end{stepbox}

\begin{notebox}[一张图可以属于多个项目]
  一张图片可以同时出现在多个项目相册里。\textbf{「加入项目相册」和「标记采用」是两件独立的事}：
  加进相册只是让它出现在项目里，并不等于采用；取消采用后，图片仍然\textbf{留在相册中}，方便你日后重新标记。
\end{notebox}

% =====================================================================
\chapter{工时与统计}
% =====================================================================

\begin{introduction}
  \item 怎么确认成员的工时
  \item 怎么导出被引张数和工时
\end{introduction}

\section{确认工时}

校区负责人在\menu{工时记录}里填报自己参与任务的工时，你负责\textbf{确认}或\textbf{退回}。

\begin{itemize}
  \item 填报人来自各校区\textbf{通讯录}；历史工时会保存当时的姓名和学号快照。
  \item 确认前系统会校验状态和版本，避免\textbf{重复确认}或并发冲突。
  \item 统计只认\textbf{已确认}的工时。
\end{itemize}

\section{导出统计}

\begin{stepbox}[导出成员数据]
  \begin{enumerate}
    \item 在\menu{数据统计}或\menu{工时记录}里选择时间范围。
    \item 点导出。系统会\textbf{在后台生成 Excel}，生成好后自动下载。
    \item 稍等片刻即可（导出是异步任务，按钮会等一会儿）。
  \end{enumerate}
\end{stepbox}

\begin{notebox}[口径说明]
  工时按\textbf{工作日期}落在所选区间来算；被引张数按项目\textbf{完成时间}落在区间、且只统计已完成项目，
  同一张图属于多个项目时\textbf{只算一次}，不会重复计发。系统\textbf{不计算薪资}，只出张数和工时。
\end{notebox}

% =====================================================================
\chapter{通讯录与好图精选}
% =====================================================================

\section{通讯录（只读）}

\menu{通讯录}里能看到各校区的成员名单。作为部长，你是\textbf{只读}的——看到的是各校区启用成员\textbf{按学号去重}
后的汇总名单。维护名单是校区负责人和管理员的事。

\begin{tipbox}[为什么拍摄者要从通讯录选]
  上传图片时，「拍摄者」必须\textbf{从通讯录里选}，不能手打。这样统计时才能准确对应到人。
  如果某人不在名单里，需要先由校区负责人把他加进通讯录。
\end{tipbox}

\section{好图精选}

\menu{好图精选}用来评选和展示部门的优秀作品。你可以发布精选、手动截止，把好图集中呈现出来。

% =====================================================================
\chapter{成员招募}
% =====================================================================

\begin{introduction}
  \item 怎么做一张报名表
  \item 怎么发布并收报名
  \item 有哪些要特别小心的地方
\end{introduction}

\section{招募是怎么运作的}

你可以设计一张\textbf{报名表}并限时发布，\textbf{同学不用登录}就能在公开页填表、上传原图作品；
提交上来的报名，你在工作站里集中查看。

\section{创建到发布}

\begin{stepbox}[发布一次招募]
  \begin{enumerate}
    \item 进入\menu{成员招募}，新建 → 设计报名表（题目、学号项、上传项）→ 保存草稿。
    \item 检查报名表无误后\textbf{发布}，并设置\textbf{开始和截止时间}。
    \item 只有状态为「已发布」且当前时间在起止区间内的招募，才会出现在公开页。
    \item 公开页地址是 \code{https://你的站点/\#/recruitment}（\code{\#} 不能少）。
  \end{enumerate}
\end{stepbox}

\begin{warnbox}[发布后表单结构会被冻结]
  为了不影响正在填表的同学，招募\textbf{一旦发布，题目、学号项、上传项就不能再改了}，只能再改名称、说明和时间。
  所以\textbf{发布前一定要检查清楚表单}。关闭招募会立即停止收新报名，已提交的不受影响。
\end{warnbox}

\section{查看报名}

报名会集中在招募详情里，用纯文本方式展示。同学上传的作品\textbf{按原样保存}，你可以预览和下载。

\begin{notebox}[学号不是实名认证]
  同一个招募里，同一个学号只能报一次（防重复报名）。但这\textbf{只是输入约束，不是身份认证}，
  不要把它当作实名信息来用。
\end{notebox}

% =====================================================================
\chapter{文档中心与消息}
% =====================================================================

\section{文档中心}

\menu{文档中心}可以编写、整理、发布部门文档（比如拍摄规范、投稿须知）。你可以拖拽整理目录结构，
并指定某篇文档是否\textbf{需要登录才能看}。

\section{消息中心}

右上角\textbf{铃铛}会提示未读消息，\menu{消息中心}里能看到所有站内通知（比如你发布需求、有人交付等）。
未读数量会在铃铛上显示角标。

\section{个人设置}

点右上角\textbf{头像}可以设置个人头像、退出登录。想改密码也从这里进入。

\begin{tipbox}[遇到问题先看这几点]
  \begin{itemize}
    \item 菜单里没有某个功能？可能是权限没开，找管理员确认你的权限组。
    \item 图片打不开 / 链接失效？多半是安全链接过期了，刷新或重新点开即可。
    \item 发布需求校区收不到通知？确认对方校区\textbf{已启用}、对方账号\textbf{已启用}。
  \end{itemize}
\end{tipbox}

\end{document}
